把学习率从 \(0.01\) 改成 \(0.011\)，训练就从收敛变成了发散。

别的什么都没动。把 float32 换成 float64 重跑，发散照样发生，所以这不是数值精度的问题。变化的不是误差的大小，而是系统的定性行为：参数更新 \(\theta_{k+1} = \theta_k - \eta\nabla L(\theta_k)\) 本身就是一个离散动力系统，学习率是它的一个参数，在某个临界值处，那个原本把轨迹吸进去的不动点变成了把轨迹推出去的。临界值算得出来——对二次目标，条件是 \(\eta < 2/\lambda_{\max}\)。

这就是本部分处理的那类问题。手上握着的是局部的变化规律，也就是每一点上``下一步往哪走''；要问的却是整体：解存不存在，能活多久，长期去向哪里，以及当参数被轻轻动一下时，这个``哪里''会不会整个换掉。

这个方向之所以重要，是因为自然与工程中的规律几乎总是先以变化率的形式出现。种群的增长率正比于当前规模，摆的角加速度由当前角度决定，膜电位的变化由当前电流决定，参数的更新量由当前梯度决定。能被直接观察或直接假设的，永远是局部的``怎样变''；而人要的答案，永远是整体的``最后成什么样''。这两者之间隔着的那段距离，就是这一部分的全部内容。

\begin{center}
\begin{tikzpicture}[x=1cm,y=1cm,font=\small,>=stealth]
  % --- 面板 A：局部规律（斜率场） ---
  \draw[black!30] (0,0) rectangle (3,2.2);
  \foreach \x in {0.4,1.0,1.6,2.2,2.6} {
    \foreach \y in {0.35,0.8,1.25,1.7} {
      \pgfmathsetmacro{\s}{1.5*\y*(2.0-\y)/2}
      \draw[black!60] (\x-0.18,\y-0.18*\s) -- (\x+0.18,\y+0.18*\s);
    }
  }
  \draw[black!85,line width=0.9pt] plot[smooth] coordinates
    {(0.15,0.25) (0.7,0.42) (1.2,0.72) (1.7,1.12) (2.2,1.5) (2.7,1.75)};
  \node[font=\footnotesize,black!70] at (1.5,-0.42) {局部：每点规定怎样变};

  % --- 面板 B：整体行为（相平面） ---
  \draw[black!30] (4.6,0) rectangle (7.4,2.2);
  \begin{scope}[shift={(6.0,1.1)}]
    \draw[black!35] (-1.25,0) -- (1.25,0);
    \draw[black!35] (0,-1.0) -- (0,1.0);
    \draw[black!85,line width=0.9pt] plot[smooth] coordinates
      {(1.2,0) (0.612,0.771) (-0.183,0.786) (-0.598,0.283)
       (-0.486,-0.240) (-0.094,-0.435) (0.232,-0.282) (0.299,0.005) (0.149,0.194)};
    \begin{scope}[rotate=155]
      \draw[black!55,line width=0.7pt] plot[smooth] coordinates
        {(1.2,0) (0.612,0.771) (-0.183,0.786) (-0.598,0.283)
         (-0.486,-0.240) (-0.094,-0.435) (0.232,-0.282) (0.299,0.005) (0.149,0.194)};
    \end{scope}
    \fill[black] (0,0) circle (0.06);
  \end{scope}
  \node[font=\footnotesize,black!70] at (6.0,-0.42) {整体：轨迹去向哪里};

  % --- 面板 C：未知量变成场 ---
  \begin{scope}[shift={(8.6,0)}]
    \draw[black!30] (0,0) rectangle (3,2.2);
    \draw[black!35] (0.1,0.2) -- (2.9,0.2);
    \draw[black!85,line width=0.9pt] plot[domain=0.1:2.9,samples=60]
      (\x,{0.2+1.4*exp(-4*(\x-1.5)*(\x-1.5))});
    \draw[black!60,line width=0.7pt] plot[domain=0.1:2.9,samples=60]
      (\x,{0.2+0.9*exp(-1.6*(\x-1.5)*(\x-1.5))});
    \draw[black!40,line width=0.7pt] plot[domain=0.1:2.9,samples=60]
      (\x,{0.2+0.5*exp(-0.7*(\x-1.5)*(\x-1.5))});
    \node[font=\footnotesize,black!70] at (1.5,-0.42) {未知量从点变成场};
  \end{scope}

  \draw[->,black!55] (3.15,1.1) -- (4.45,1.1);
  \draw[->,black!55] (7.55,1.1) -- (8.45,1.1);
\end{tikzpicture}

\emph{三块面板问的是同一个问题，只是未知量一级级升维：先是一个数，再是相空间里的一条轨迹，最后是空间上的一整个场。每升一级，``局部决定整体''这句话就需要一套新的语言来兑现。}
\end{center}

\subsection{五条贯穿始终的判断}

\textbf{局部规律不保证整体解存在。}\(y' = y^2\)、\(y(0)=1\) 的解是 \(1/(1-t)\)，在 \(t=1\) 处炸掉——右端处处光滑，解却只活了有限的一段时间。而 \(y' = \sqrt{y}\)、\(y(0)=0\) 有无穷多个解，因为右端在原点不 Lipschitz。这两个反例挡住的是两件不同的事：Lipschitz 条件管的是唯一性，右端的增长速度管的是解能活多久。分开记账，才知道手上这个方程缺的到底是哪一个。

\textbf{解不出来不等于看不懂。}线性系统可以完全解出，矩阵指数 \(e^{At}\) 把答案写成闭式；非线性系统绝大多数解不出来。但相平面上的不动点、轨道的走向、Lyapunov 函数给出的能量下降，全都是不依赖解析解的定性结论——而定性结论往往才是真正要的那个。问一个控制系统稳不稳，比问它的精确轨迹有用得多。这一部分把大量篇幅放在几何与能量上而不是求解技巧上，理由就在这里。

\textbf{定性行为可以对参数不连续地依赖。}参数连续地变，不动点的数目与稳定性却会突变，这就是分岔。开头那个学习率的例子是最朴素的一个。``我只改了百分之十''这种直觉在动力系统里不成立，因为决定行为的不是参数动了多少，而是它有没有越过某个临界值；而临界值的位置通常与系统的谱有关。这条线一直通到深度学习的训练稳定性和控制里的稳定裕度。

\textbf{确定不等于可预测。}混沌系统完全确定，没有任何随机成分，但初值的微小差别会被指数放大，越过某个时间视界之后预测就失去意义。Lyapunov 指数量化的正是这个放大率，它回答的是``预测能维持多久''，不是``能不能预测''。把这两个词分开很要紧：混淆它们，就会对长期模拟的结果产生不该有的信任。

\textbf{写下一个 PDE 不等于把问题提好了。}还要说清边界条件、初始数据，以及``解''这个词指的是什么。适定性的三条里最容易塌的是第三条——对数据的连续依赖。热方程往回走就是反问题，输入的微小噪声被 \(e^{\kappa k^2 t}\) 放大，条件数没有上界，所以图像去模糊必须靠正则化，而正则化换来的代价是分辨率。再往前一步，激波处的解根本不可微，逐点形式的方程已经不成立，只能换成积分形式：解的概念本身是需要声明的。

\subsection{三章怎么安排}

\hyperref[ch:134]{``常微分方程理论：解是否存在、是否唯一''}是地基。它先用增长、衰减与振动这三种最简单的行为把``方程即变化规律''坐实，再以线性系统作为可以完全分析的样板——矩阵指数在这里出现，特征值的实部决定长期趋势，虚部决定振荡，这一条后面反复要用。存在唯一性那一篇建议读完整，因为它的价值不在定理本身，而在两个反例分别逼出了哪一个条件。最后一篇讲尺度分离与渐近近似，处理刚性系统与多时间尺度，与数值方法那边的步长限制是同一回事。

\hyperref[ch:135]{``稳定性与动力系统：长期行为的几何与能量''}换一套眼光看同样的方程。相平面把解画成几何对象，于是``长期去哪里''变成了看图就能回答的问题；Lyapunov 函数则把稳定性判断从求解降级成找一个一直在下降的量——这个思路在强化学习的收敛证明与优化算法的分析里会再次出现。最后一篇讲离散迭代、分岔与混沌，是全书里离``为什么调个参数就变天''最近的一篇。

\hyperref[ch:136]{``偏微分方程初步：场随时间与空间一起变化''}把未知量从有限维的状态升级成空间上的场。重点不是求解技巧，是三件事：方程的类型决定行为（抛物抹平、双曲搬运、椭圆由边界锁死），适定性要逐条对账，以及解的概念在什么时候必须放宽。分离变量那一篇顺带说清了 Fourier 展开为什么在最对称的几何上才好用；最后一篇引入弱解与半群，是给需要严格框架的读者准备的。

这一部分与\hyperref[ch:066]{``常微分方程数值方法：沿局部斜率走完整条轨迹''}是有意分工的。那里回答``怎么把解算准''，这里回答排在它前面的问题：解是否存在、是否唯一、长期是什么形状。顺序不能颠倒——若解根本不存在或不唯一，再精细的步长控制也只是在精确地逼近一个不存在的对象。往后看，\hyperref[ch:137]{``控制与强化学习：让系统按我们的意愿演化''}接手的是下一个问题：这里只是旁观系统自己往哪走，那里要往里面加一只手，让它往我们要的地方走。稳定性的语言到那边会原样复用，只是判断对象从``它稳不稳''换成了``我能不能让它稳''。
